\begin{filecontents}{refs.bib}
@article{Ahn2025A,
title={A Novel Mixed Integer Programming Model With Precedence Relation-Based Decision Variables for Non-Cyclic Scheduling of Cluster Tools},
author={Jeongsun Ahn and Hyun-Jung Kim},
journal={IEEE Transactions on Automation Science and Engineering},
year={2025},
volume={22},
pages={2893-2908},
doi={10.1109/tase.2024.3386562}
}

@article{Al-Khatib2025A,
title={A comparative analysis of MIP formulations for the multi-mode resource-constrained independent task scheduling problem},
author={Maryam Al-Khatib and M. Haouari and M. Kharbeche},
journal={Operational Research},
year={2025},
volume={25},
doi={10.1007/s12351-025-00924-2}
}

@article{Ansótegui2019Automatic,
title={Automatic Detection of At-Most-One and Exactly-One Relations for Improved SAT Encodings of Pseudo-Boolean Constraints},
author={C. Ansótegui and Miquel Bofill and Jordi Coll and Nguyen Dang and J. Esteban and Ian Miguel and Peter William Nightingale and András Z. Salamon and Josep Suy and Mateu Villaret},
year={2019},
pages={20-36},
doi={10.1007/978-3-030-30048-7_2}
}

@article{Bofill2022The,
title={The Sample Analysis Machine Scheduling Problem: Definition and comparison of exact solving approaches},
author={Miquel Bofill and Jordi Coll and G. Martín and Josep Suy and Mateu Villaret},
journal={Comput. Oper. Res.},
year={2022},
volume={142},
pages={105730},
doi={10.1016/j.cor.2022.105730}
}

@article{Bofill2020SMT,
title={SMT encodings for Resource-Constrained Project Scheduling Problems},
author={Miquel Bofill and Jordi Coll and Josep Suy and Mateu Villaret},
journal={Comput. Ind. Eng.},
year={2020},
volume={149},
pages={106777},
doi={10.1016/j.cie.2020.106777}
}

@article{Bofill2021SAT,
title={SAT Encodings for Pseudo-Boolean Constraints Together With At-Most-One Constraints},
author={Miquel Bofill and Jordi Coll and Peter William Nightingale and Josep Suy and Felix Ulrich-Oltean and Mateu Villaret},
journal={Artif. Intell.},
year={2021},
volume={302},
pages={103604},
doi={10.1016/j.artint.2021.103604}
}

@article{Bofill2020An,
title={An MDD-based SAT encoding for pseudo-Boolean constraints with at-most-one relations},
author={Miquel Bofill and Jordi Coll and Josep Suy and Mateu Villaret},
journal={Artificial Intelligence Review},
year={2020},
volume={53},
pages={5157 - 5188},
doi={10.1007/s10462-020-09817-6}
}

@article{Bofill2023SAT,
title={SAT Encodings for Pseudo-Boolean Constraints Together With At-Most-One Constraints (Extended Abstract)},
author={Miquel Bofill and Jordi Coll and Peter William Nightingale and Josep Suy and Felix Ulrich-Oltean and Mateu Villaret},
year={2023},
pages={6853-6857},
doi={10.24963/ijcai.2023/769}
}

@article{Caselli2023Exact,
title={Exact algorithms for a parallel machine scheduling problem with workforce and contiguity constraints},
author={Giulia Caselli and Maxence Delorme and M. Iori and C. Magni},
journal={Comput. Oper. Res.},
year={2023},
volume={163},
pages={106484},
doi={10.1016/j.cor.2023.106484}
}

@article{DeAzevedo2021A,
title={A satisfiability and workload-based exact method for the resource constrained project scheduling problem with generalized precedence constraints},
author={Guilherme Henrique Ismael de Azevedo and A. Pessoa and A. Subramanian},
journal={Eur. J. Oper. Res.},
year={2021},
volume={289},
pages={809-824},
doi={10.1016/j.ejor.2019.07.056}
}

@article{Edwards2019Symmetry,
title={Symmetry breaking of identical projects in the high-multiplicity RCPSP/max},
author={Steven J. Edwards and Davaatseren Baatar and K. Smith‐Miles and Andreas T. Ernst},
journal={Journal of the Operational Research Society},
year={2019},
volume={72},
pages={1822 - 1843},
doi={10.1080/01605682.2019.1595192}
}

@article{Fan2021Genetic,
title={Genetic programming-based hyper-heuristic approach for solving dynamic job shop scheduling problem with extended technical precedence constraints},
author={Huali Fan and Hegen Xiong and M. Goh},
journal={Comput. Oper. Res.},
year={2021},
volume={134},
pages={105401},
doi={10.1016/j.cor.2021.105401}
}

@article{Hölldobler2013An,
title={An Efficient Encoding of the at-most-one Constraint},
author={Steffen Hölldobler and Van-Hau Nguyen},
year={2013}
}

@article{Kasapidis2023On,
title={On the multiresource flexible job‐shop scheduling problem with arbitrary precedence graphs},
author={Gregory A. Kasapidis and S. Dauzére-Pérés and D. Paraskevopoulos and Panagiotis P. Repoussis and C. Tarantilis},
journal={Production and Operations Management},
year={2023},
volume={32},
pages={2322 - 2330},
doi={10.1111/poms.13977}
}

@article{Klein2024Mixed-integer,
title={Mixed-integer linear programming for project scheduling under various resource constraints},
author={Nicklas Klein and Mario Gnägi and Norbert Trautmann},
journal={Eur. J. Oper. Res.},
year={2024},
volume={319},
pages={79-88},
doi={10.1016/j.ejor.2024.06.036}
}

@article{Koné2009Event-based,
title={Event-based MILP models for resource-constrained project scheduling problems},
author={Oumar Koné and C. Artigues and P. Lopez and M. Mongeau},
journal={Comput. Oper. Res.},
year={2009},
volume={38},
pages={3-13},
doi={10.1016/j.cor.2009.12.011}
}

@article{Lenstra1978Complexity,
title={Complexity of Scheduling under Precedence Constraints},
author={J. Lenstra and A. Kan},
journal={Oper. Res.},
year={1978},
volume={26},
pages={22-35},
doi={10.1287/opre.26.1.22}
}

@article{Meng2023Transitive,
title={Transitive Reduction Approach to Large-Scale Parallel Machine Rescheduling Problem With Controllable Processing Times, Precedence Constraints and Random Machine Breakdown},
author={De Meng and Zu-ren Feng and Zhi Ren and Yue Liu},
journal={IEEE Access},
year={2023},
volume={11},
pages={7727-7738},
doi={10.1109/access.2023.3238639}
}

@article{Naderi2023Mixed-Integer,
title={Mixed-Integer Programming vs. Constraint Programming for Shop Scheduling Problems: New Results and Outlook},
author={B. Naderi and Rubén Ruiz and V. Roshanaei},
journal={INFORMS J. Comput.},
year={2023},
volume={35},
pages={817-843},
doi={10.1287/ijoc.2023.1287}
}

@article{Prot2017A,
title={A survey on how the structure of precedence constraints may change the complexity class of scheduling problems},
author={Damien Prot and Odile Bellenguez-Morineau},
journal={Journal of Scheduling},
year={2017},
volume={21},
pages={3 - 16},
doi={10.1007/s10951-017-0519-z}
}

@article{Sayah2023Continuous-Time,
title={Continuous-Time Formulations for Multi-Mode Project Scheduling},
author={D. Sayah},
journal={Comput. Oper. Res.},
year={2023},
volume={152},
pages={106147},
doi={10.1016/j.cor.2023.106147}
}

@article{Schnell2015On,
title={On the efficient modeling and solution of the multi-mode resource-constrained project scheduling problem with generalized precedence relations},
author={Alexander Schnell and R. Hartl},
journal={OR Spectrum},
year={2015},
volume={38},
pages={283 - 303},
doi={10.1007/s00291-015-0419-6}
}

@article{Tesch2017Improved,
title={Improved Compact Models for the Resource-Constrained Project Scheduling Problem},
author={Alexander Tesch},
year={2017},
pages={25-30},
doi={10.1007/978-3-319-55702-1_4}
}

@article{Tesch2020A,
title={A polyhedral study of event-based models for the resource-constrained project scheduling problem},
author={Alexander Tesch},
journal={Journal of Scheduling},
year={2020},
volume={23},
pages={233-251},
doi={10.1007/s10951-020-00647-6}
}

@article{Truong2025Sequential,
title={Sequential Counter Encoding for Staircase At-Most-One Constraints},
author={Hieu Xuan Truong and Tuyen Van Kieu and Khanh Van To},
year={2025},
pages={164-175},
doi={10.5220/0013124600003890}
}
\end{filecontents}

\documentclass{article}
\usepackage{adjustbox}
\begin{document}

\section*{\textbf{SAMSP}: hướng cải tiến \textbf{compact SMT/MIP}, \textbf{staircase AMO}, và \textbf{extended precedence}}


Bạn có đủ cơ sở tài liệu để viết một bản thảo mạnh: hướng cải tiến hợp lý nhất cho \textit{The Sample Analysis Machine Scheduling Problem} là giảm phụ thuộc theo chân trời thời gian của mô hình SMT/MILP, tận dụng cấu trúc AMO/PB để nén số biến–số mệnh đề, và khai thác mạnh hơn đồ thị tiền tố mở rộng để vừa cắt giảm miền biến vừa tăng suy diễn tiền xử lý.


\subsection*{Bài Toán Gốc Và Khoảng Trống}


SAMSP được Bofill và cộng sự mô hình hóa như một trường hợp đặc biệt của RCPSP-Cu, trong đó hoạt động vừa chịu ràng buộc tiền tố, tài nguyên tái tạo, vừa có chiếm chỗ/trả chỗ trên các vùng lưu trữ hữu hạn  \cite{Bofill2022The}. Điểm khó riêng của SAMSP là chiếm chỗ bắt đầu ngay khi hoạt động bắt đầu, còn giải phóng chỗ chỉ xảy ra khi hoạt động kết thúc, khiến nó không rút gọn trực tiếp về RCPSP chuẩn với thời lượng cố định và tài nguyên tái tạo thuần túy  \cite{Bofill2022The}.


Bài gốc so sánh CP, SMT và MILP trên các instance công nghiệp thực, và kết luận CP là tốt nhất, SMT đứng thứ hai, còn MILP tụt khá xa  \cite{Bofill2022The}. Nguyên nhân quan trọng là SMT và MILP của bài gốc đều là mô hình time-indexed, nên số biến tăng theo cận trên makespan; chính bài báo nhấn mạnh rằng kích thước SMT tỉ lệ với upper bound của lịch  \cite{Bofill2022The}. Điều này đặc biệt bất lợi cho SAMSP vì các instance có makespan lớn do vùng observation là nút cổ chai khan hiếm  \cite{Bofill2022The}.


\subsection*{Bằng Chứng Cho Hướng Compact}


Về tổng quan, tài liệu RCPSP cho thấy các mô hình compact có lợi thế vì kích thước mạnh theo số công việc thay vì phụ thuộc trực tiếp vào chân trời thời gian  \cite{Tesch2017Improved} \cite{Koné2009Event-based}. Koné và cộng sự chỉ ra rõ rằng các event-based MILP cho RCPSP có ít biến hơn time-indexed và không phụ thuộc vào scheduling horizon, nên phù hợp hơn khi horizon lớn  \cite{Koné2009Event-based}.


Bằng chứng so sánh cũng khá nhất quán rằng không có một MIP formulation thống trị mọi loại instance, nhưng khi horizon lớn thì event-based hoặc continuous-time thường đáng thử hơn time-indexed  \cite{Koné2009Event-based}. Tesch còn cho thấy trong họ mô hình event-based cho RCPSP, có thể sắp hạng polyhedral theo thứ tự IEE ≻ SEE ≻ OOE, đồng thời nhiều phép biến đổi giữa các mô hình là tuyến tính, nên cải tiến cấu trúc formulation hoàn toàn có thể được biện minh lý thuyết chứ không chỉ thực nghiệm  \cite{Tesch2020A}.


Bài toán gần với SAMSP hơn là RCPSP-CPR cũng cho thấy mô hình compact continuous-time với biến sequencing mới có thể vượt các mô hình MILP chuẩn trên cả RCPSP cổ điển lẫn tài nguyên kiểu storage/consumption-production  \cite{Klein2024Mixed-integer}. Với bài toán cluster tool phi chu kỳ, Ahn và Kim cũng giảm số biến bằng cách chuyển từ position-based variables sang precedence-relation-based variables, và báo cáo thời gian tính tốt hơn mô hình cũ  \cite{Ahn2025A}.


\subsection*{So Sánh Các Họ Mô Hình}


Các họ mô hình liên quan trực tiếp đến hướng bạn đề xuất được tóm tắt dưới đây.


\mbox{}\\
\begin{adjustbox}{max width=\textwidth}
\begin{tabular}{p{4.0cm} | p{4.0cm} | p{4.0cm} | p{4.0cm}}
\hline
Họ mô hình & Ý tưởng nén chính & Điểm mạnh cho SAMSP & Rủi ro chính \\
\hline
Time-indexed DT/DDT & Biến theo hoạt động × thời điểm & LP relaxation mạnh trên horizon ngắn  \cite{Koné2009Event-based} & Bùng nổ theo makespan  \cite{Koné2009Event-based} \\
Event-based MILP & Biến theo event, không theo thời gian & Phù hợp horizon lớn  \cite{Koné2009Event-based} & Relaxation có thể yếu hơn DDT ở benchmark truyền thống  \cite{Koné2009Event-based} \\
Continuous-time/compact sequencing & Biến thứ tự hoặc song song cần thiết & Tự nhiên cho precedence và storage sequencing  \cite{Klein2024Mixed-integer} & Có thể cần big-M hoặc tách trường hợp tinh tế  \cite{Koné2009Event-based} \cite{Sayah2023Continuous-Time} \\
SMT/SAT với PB(AMO) & Mã hóa tài nguyên bằng SAT compact & Giảm mạnh số biến/mệnh đề và tăng tốc  \cite{Bofill2020SMT} \cite{Bofill2021SAT} \\
\hline
\end{tabular}
\end{adjustbox}


\textbf{Figure 1:} Các họ formulation phù hợp để cải tiến SAMSP


Một điểm quan trọng cho phần positioning là kết quả hiện đại không còn xem MIP mặc định tốt hơn CP. Trên 12 họ bài toán scheduling, CP đặt ra giới hạn mới về kích thước bài toán có thể giải bằng exact off-the-shelf methods  \cite{Naderi2023Mixed-Integer}. Vì vậy, nếu bản thảo của bạn đề xuất compact SMT/MIP thay vì cố “đánh bại CP” trên mọi instance, lập luận sẽ thuyết phục hơn: mục tiêu nên là thu hẹp khoảng cách với CP trong các lớp instance mà time-indexed SMT/MILP gốc bị phạt bởi horizon lớn.


\subsection*{Staircase, Ladder Và PB(AMO)}


Trong bài gốc, biến thể SMT tốt nhất trên tài nguyên unary dùng AMO với regular/ladder encoding cho renewable resources  \cite{Bofill2022The}. Đây là điểm vào rất tự nhiên cho đóng góp staircase. Về nền tảng, AMO là ràng buộc cốt lõi trong SAT encoding, và các encoding ladder, regular, sequential có quan hệ gần nhau; các biến thể mới như bimander có thể cạnh tranh mạnh nhưng lựa chọn tốt nhất phụ thuộc cấu trúc bài toán  \cite{Hölldobler2013An}.


Dòng nghiên cứu PB(AMO) của chính nhóm tác giả cho thấy khi ràng buộc pseudo-Boolean đi kèm các AMO trên tập con biến, encoding có thể nhỏ hơn đáng kể và thời gian giải cải thiện hơn một bậc độ lớn trong một số trường hợp  \cite{Bofill2021SAT}. Ý tưởng cốt lõi là không biểu diễn những gán đã bị AMO cấm, nhờ đó decision diagram hoặc totalizer nhỏ hơn  \cite{Bofill2020An}. Hơn nữa, phát hiện tự động AMO/EO cũng đã được chứng minh giúp giảm mạnh kích thước encoding và cải thiện thời gian giải  \cite{Ansótegui2019Automatic}.


Bài mới về Staircase AMO còn gần đề xuất của bạn hơn: các AMO liên tiếp có cấu trúc staircase do chia sẻ mạnh giữa các block kế cận, và biểu diễn bằng sequential counter theo staircase cho ít biến và ít mệnh đề hơn so với BDD-based SCAMO, đồng thời cho thời gian giải tốt hơn  \cite{Truong2025Sequential}. Điều này tạo ra một luận điểm mới cho bản thảo: nếu các precedence/resource conflicts của SAMSP sinh ra họ AMO chồng lấn theo cửa sổ thời gian hoặc theo extended precedence windows, thì staircase sharing không chỉ là thay encoding cục bộ mà là khai thác cấu trúc lặp của toàn bộ họ mệnh đề.


\subsection*{Extended Precedence, Closure Và Reduction}


Bài gốc đã dùng transitive closure của precedence graph để tạo extended precedence graph, suy ra lower/upper bounds và time windows, rồi thêm symmetry-breaking precedences trước khi tính closure để giảm miền biến  \cite{Bofill2022The}. Kết quả thực nghiệm cho thấy symmetry breaking làm tăng mạnh số precedence mở rộng nhưng đồng thời cải thiện lower bound và giảm đáng kể thời gian giải cho cả CP, SMT và MILP  \cite{Bofill2022The}.


\mbox{}\\
\begin{adjustbox}{max width=\textwidth}
\begin{tabular}{p{5.3cm} | p{5.3cm} | p{5.3cm}}
\hline
Kỹ thuật trên precedence graph & Tác dụng được báo cáo & Hàm ý cho bản thảo \\
\hline
Transitive closure / extended precedence & Thu hẹp cửa sổ thời gian, sinh implied constraints  \cite{Bofill2022The} & Nên là baseline bắt buộc \\
Symmetry-breaking by added precedences & Tăng LB và giảm runtime lớn  \cite{Bofill2022The} & Có thể kết hợp với compact encoding \\
Transitive reduction & Loại precedence dư thừa, giảm thời gian và RAM  \cite{Meng2023Transitive} & Nên thử song song với closure \\
Generalized/extended precedence modeling & Tăng sức mạnh exact methods khi có propagator riêng  \cite{Schnell2015On} & Hỗ trợ lập luận cho precedence-aware encoding \\
\hline
\end{tabular}
\end{adjustbox}


\textbf{Figure 2:} Các thao tác chính trên đồ thị precedence


Điểm đáng chú ý là closure và reduction không mâu thuẫn; chúng phục vụ hai tầng khác nhau. Closure giúp suy diễn implied precedences và cắt miền; reduction giúp loại bớt cung dư thừa trong mô hình hoặc solver input  \cite{Meng2023Transitive}. Với bản thảo của bạn, đây là một đóng góp tốt: thử nghiệm riêng các chế độ \textit{raw graph}, \textit{transitive reduction}, \textit{extended closure}, và \textit{reduction + closure-derived windows}. Tesch còn chỉ ra rằng chỉ riêng tích hợp precedence transitive information đã loại được nhiều gán event không khả thi trong event-based models  \cite{Tesch2020A}.


Ngoài RCPSP cổ điển, tài liệu rộng hơn cũng xác nhận cấu trúc precedence quyết định mạnh độ khó tính toán; nhiều bài toán trở nên NP-hard dưới precedence tổng quát nhưng dễ hơn trên cấu trúc đặc biệt  \cite{Prot2017A} \cite{Lenstra1978Complexity}. Các mô hình với arbitrary precedence graphs, nonlinear routes, hoặc ETPC đều nhấn mạnh rằng mở rộng precedence là thực tế công nghiệp chứ không chỉ là biến thể lý thuyết  \cite{Kasapidis2023On} \cite{Fan2021Genetic}.


\subsection*{Cách Định Vị Đóng Góp Mới}


Một framing thuyết phục cho bài của bạn là: bài gốc đã chứng minh SAMSP là RCPSP-Cu công nghiệp khó, nhưng phần SMT/MILP còn bị giới hạn bởi time indexing và chưa khai thác hết cấu trúc AMO/PB cũng như cấu trúc precedence mở rộng  \cite{Bofill2022The}. Đóng góp mới vì thế có thể được phát biểu thành ba tầng.


\begin{itemize}
\item \textbf{Tầng mô hình}: đề xuất compact formulation giảm phụ thuộc vào horizon, theo tinh thần event-based/continuous-time/predecessence-variable formulations  \cite{Koné2009Event-based} \cite{Klein2024Mixed-integer} \cite{Ahn2025A}.
\item \textbf{Tầng mã hóa SAT/SMT}: thay AMO rời rạc bằng staircase-aware PB(AMO)/SCAMO encoding để giảm số biến và số clause  \cite{Bofill2021SAT} \cite{Truong2025Sequential} \cite{Bofill2023SAT}.
\item \textbf{Tầng tiền xử lý precedence}: so sánh closure, symmetry breaking, transitive reduction, và extended precedence windows như cơ chế giảm biến trước khi solve  \cite{Bofill2022The} \cite{Edwards2019Symmetry} \cite{Meng2023Transitive}.
\end{itemize}


Một bảng thực nghiệm nên tách ít nhất bốn chỉ số: số biến, số mệnh đề/ràng buộc, thời gian encode, thời gian solve; và thêm chất lượng LB/UB nếu bạn giữ MILP hoặc SMT tối ưu hóa từng bước. Điều này phù hợp với truyền thống so sánh formulation trong RCPSP và scheduling nói chung, nơi kích thước mô hình và độ mạnh relaxation đều được xem là đầu ra khoa học chính, không chỉ runtime  \cite{Koné2009Event-based} \cite{Tesch2017Improved} \cite{Al-Khatib2025A}.


\subsection*{Khoảng Trống Nghiên Cứu Và Đề Cương Viết}


Khoảng trống rõ nhất trong văn liệu là chưa có công trình nào, trong tập tài liệu này, kết hợp trực tiếp ba ý tưởng cho một bài toán kiểu SAMSP: storage-aware RCPSP-Cu, compact exact formulation, và staircase PB(AMO) encoding dựa trên extended precedence. Các nhánh này hiện tồn tại rời rạc: RCPSP-Cu/storage compact MILP  \cite{Klein2024Mixed-integer}, SMT/PB(AMO) cho scheduling  \cite{Bofill2020SMT} \cite{Bofill2021SAT}, và arbitrary/generalized precedence exact solving  \cite{DeAzevedo2021A} \cite{Schnell2015On}.


Câu hỏi mở lớn nhất là liệu giảm số biến và số mệnh đề có chuyển thành giảm runtime trên chính các instance SAMSP hay không. Văn liệu cho thấy điều này thường đúng nhưng không tự động: formulation nhỏ hơn chưa chắc relaxation mạnh hơn, và không có formulation nào thống trị mọi họ instance  \cite{Koné2009Event-based} \cite{Caselli2023Exact}. Vì vậy, bản thảo nên tránh tuyên bố “compact hơn nên tốt hơn”, mà nên kiểm định theo lớp instance, đặc biệt các lớp có makespan lớn, observation bottleneck mạnh, và nhiều test cùng loại nơi symmetry/precedence structure rõ nhất  \cite{Bofill2022The}.


Dịch chuyển quan trọng nhất mà bằng chứng hiện có gợi ra là: trong các bài toán kiểu SAMSP, hiệu năng exact solving không chỉ do solver quyết định mà chủ yếu do mức độ khai thác cấu trúc của precedence, AMO và storage trong formulation. Câu hỏi mở lớn nhất cho bản thảo của bạn là cấu trúc staircase trên ladder constraints có còn giữ lợi thế khi ghép với extended precedence graph và cumulative/storage constraints của SAMSP ở quy mô công nghiệp hay không.

% These search results were found and analyzed using Consensus, an AI-powered search engine for research. Try it at https://consensus.app. © 2026 Consensus NLP, Inc. Personal, non-commercial use only; redistribution requires copyright holders’ consent.

\bibliographystyle{plain}
\bibliography{refs}
\end{document}